% GENERATED FILE. Edit canonical lessons, then run npm run generate:books.
% canonical-source-hash: fnv1a64:81b9052a6c8abf96
% canonical-lessons: AR-C296-sahafa, AR-C296-muharrir, AR-C296-maqal, AR-C296-harrara, AR-C296-nashra

\chapter{The press}
\label{ch:ar-the-press}
\hlchaptermodality{296}

\begin{chapteropening}
\textbf{By the end of this chapter:} \emph{I can say the press, an editor, an article, to edit and a bulletin.}

Complete the last lesson of chapter 296: I can say the press, an editor, an article, to edit and a bulletin.
\end{chapteropening}

\section[\texorpdfstring{\ar{صحافة} (ṣaḥāfa)}{صحافة (ṣaḥāfa)}]{\ar{صحافة} (ṣaḥāfa) — journalism; the press}

\label{lesson:AR-C296-sahafa}



\begin{quote}
Before the new one: say the Arabic for to recount, then the Arabic for one day.
\end{quote}

\subsection*{You'll want to know: \ar{صحافة}}
\textbf{\ar{صحافة}} — \emph{ṣaḥāfa} — \textquotedblleft{}journalism; the press\textquotedblright{}.

\begin{quote}
\ar{درس} \ar{أخي} \ar{الصحافة} \ar{في} \ar{الجامعة}.
\end{quote}

\begin{quote}
\emph{darasa akhī aṣ-ṣaḥāfa fī al-jāmiʿa.}
\end{quote}

\begin{quote}
My brother studied journalism at university.
\end{quote}

A journalist is \ar{صحفي}; the newspaper is \ar{صحيفة}: one root.

\begin{cousinweb}
\textbf{\ar{صحافة}} (\emph{ṣaḥāfa}) is built on the native Arabic root \emph{ṣ-ḥ-f} (Lane), a written leaf or page, the root of \emph{ṣaḥīfa}. The modern sense, the press, comes from \emph{ṣaḥīfa}, a newspaper.
\end{cousinweb}

\begin{grammarlens}[title={What you've built}]
One more word for telling what happened.
\end{grammarlens}

\subsection*{Your turn}
\begin{itemize}
\raggedright
  \item \emph{Say it:} \emph{ṣaḥāfa}
  \item \emph{Say it:} \emph{ṣaḥāfa}, once more
  \item \emph{Say it:} \emph{darasa akhī aṣ-ṣaḥāfa fī al-jāmiʿa.} — My brother studied journalism at university.
  \item \emph{Recall:} say the Arabic for to recount, then the Arabic for one day, then say \emph{ṣaḥāfa} again
\end{itemize}

\begin{tcolorbox}[breakable,colback=teal!4,colframe=teal!35!black,title={Before you move on}]
What does \ar{صحافة} mean? (\textquotedblleft{}Journalism; the press\textquotedblright{}.) Say it once more.
\end{tcolorbox}
\section[\texorpdfstring{\ar{محرر} (muḥarrir)}{محرر (muḥarrir)}]{\ar{محرر} (muḥarrir) — an editor}

\label{lesson:AR-C296-muharrir}



\begin{quote}
Before the new one: say the Arabic for one day, then the Arabic for journalism.
\end{quote}

\subsection*{You'll want to know: \ar{محرر}}
\textbf{\ar{محرر}} — \emph{muḥarrir} — \textquotedblleft{}an editor\textquotedblright{}.

\begin{quote}
\ar{يعمل} \ar{خالي} \ar{محررا} \ar{في} \ar{مجلة}.
\end{quote}

\begin{quote}
\emph{yaʿmalu khālī muḥarriran fī majalla.}
\end{quote}

\begin{quote}
My uncle works as an editor at a magazine.
\end{quote}

The verbal noun \ar{تحرير} is editing: the editorial office is \ar{قسم} \ar{التحرير}.

\begin{cousinweb}
\textbf{\ar{محرر}} (\emph{muḥarrir}) is built on the native Arabic root \emph{ḥ-r-r} (Lane), \textquotedblleft{}to be free\textquotedblright{}, the root of \emph{ḥurr}, free. Form II \emph{ḥarrara} means \textquotedblleft{}to set free\textquotedblright{} and \textquotedblleft{}to write a text out clean\textquotedblright{}; \emph{muḥarrir} is its participle.
\end{cousinweb}

\begin{grammarlens}[title={What you've built}]
One more word for telling what happened.
\end{grammarlens}

\subsection*{Your turn}
\begin{itemize}
\raggedright
  \item \emph{Say it:} \emph{muḥarrir}
  \item \emph{Say it:} \emph{muḥarrir}, once more
  \item \emph{Say it:} \emph{yaʿmalu khālī muḥarriran fī majalla.} — My uncle works as an editor at a magazine.
  \item \emph{Recall:} say the Arabic for one day, then the Arabic for journalism, then say \emph{muḥarrir} again
\end{itemize}

\begin{tcolorbox}[breakable,colback=teal!4,colframe=teal!35!black,title={Before you move on}]
What does \ar{محرر} mean? (\textquotedblleft{}An editor\textquotedblright{}.) Say it once more.
\end{tcolorbox}
\section[\texorpdfstring{\ar{مقال} (maqāl)}{مقال (maqāl)}]{\ar{مقال} (maqāl) — an article, an essay}

\label{lesson:AR-C296-maqal}



\begin{quote}
Before the new one: say the Arabic for journalism, then the Arabic for an editor.
\end{quote}

\subsection*{You'll want to know: \ar{مقال}}
\textbf{\ar{مقال}} — \emph{maqāl} — \textquotedblleft{}an article, an essay\textquotedblright{}.

\begin{quote}
\ar{نشرت} \ar{الصحيفة} \ar{مقالا} \ar{عن} \ar{البحر}.
\end{quote}

\begin{quote}
\emph{nasharat aṣ-ṣaḥīfa maqālan ʿan al-baḥr.}
\end{quote}

\begin{quote}
The newspaper published an article about the sea.
\end{quote}

Plural \ar{مقالات}. \ar{مقالة} is the same word.

\begin{cousinweb}
\textbf{\ar{مقال}} (\emph{maqāl}) is built on the native Arabic root \emph{q-w-l} (Lane), \textquotedblleft{}to say\textquotedblright{}, the root of \emph{qāla}. \emph{Maqāl} is first \textquotedblleft{}a saying\textquotedblright{}; in the modern press it is a written piece.
\end{cousinweb}

\begin{grammarlens}[title={What you've built}]
One more word for telling what happened.
\end{grammarlens}

\subsection*{Your turn}
\begin{itemize}
\raggedright
  \item \emph{Say it:} \emph{maqāl}
  \item \emph{Say it:} \emph{maqāl}, once more
  \item \emph{Say it:} \emph{nasharat aṣ-ṣaḥīfa maqālan ʿan al-baḥr.} — The newspaper published an article about the sea.
  \item \emph{Recall:} say the Arabic for journalism, then the Arabic for an editor, then say \emph{maqāl} again
\end{itemize}

\begin{tcolorbox}[breakable,colback=teal!4,colframe=teal!35!black,title={Before you move on}]
What does \ar{مقال} mean? (\textquotedblleft{}An article, an essay\textquotedblright{}.) Say it once more.
\end{tcolorbox}
\section[\texorpdfstring{\ar{حرر} (ḥarrara)}{حرر (ḥarrara)}]{\ar{حرر} (ḥarrara) — to edit; to set free}

\label{lesson:AR-C296-harrara}



\begin{quote}
Before the new one: say the Arabic for an editor, then the Arabic for an article.
\end{quote}

\subsection*{You'll want to know: \ar{حرر}}
\textbf{\ar{حرر}} — \emph{ḥarrara} — \textquotedblleft{}to edit; to set free\textquotedblright{}.

\begin{quote}
\ar{حررت} \ar{أختي} \ar{كتابا} \ar{عن} \ar{الطبخ}.
\end{quote}

\begin{quote}
\emph{ḥarrarat ukhtī kitāban ʿan aṭ-ṭabkh.}
\end{quote}

\begin{quote}
My sister edited a book about cooking.
\end{quote}

Present: \ar{يحرر}, \emph{yuḥarrir}.

Of people and lands it means to free: \ar{حرر} \ar{الجيش} \ar{المدينة}.

\begin{cousinweb}
\textbf{\ar{حرر}} (\emph{ḥarrara}) is built on the native Arabic root \emph{ḥ-r-r} (Lane), \textquotedblleft{}to be free\textquotedblright{}. It is Form II, \textquotedblleft{}to make free\textquotedblright{}, and of a text \textquotedblleft{}to write it out clean and correct\textquotedblright{}.
\end{cousinweb}

\begin{grammarlens}[title={What you've built}]
One more word for telling what happened.
\end{grammarlens}

\subsection*{Your turn}
\begin{itemize}
\raggedright
  \item \emph{Say it:} \emph{ḥarrara}
  \item \emph{Say it:} \emph{ḥarrara}, once more
  \item \emph{Say it:} \emph{ḥarrarat ukhtī kitāban ʿan aṭ-ṭabkh.} — My sister edited a book about cooking.
  \item \emph{Recall:} say the Arabic for an editor, then the Arabic for an article, then say \emph{ḥarrara} again
\end{itemize}

\begin{tcolorbox}[breakable,colback=teal!4,colframe=teal!35!black,title={Before you move on}]
What does \ar{حرر} mean? (\textquotedblleft{}To edit; to set free\textquotedblright{}.) Say it once more.
\end{tcolorbox}
\section[\texorpdfstring{\ar{نشرة} (nashra)}{نشرة (nashra)}]{\ar{نشرة} (nashra) — a bulletin, a news broadcast}

\label{lesson:AR-C296-nashra}



\begin{quote}
Before the new one: say the Arabic for an article, then the Arabic for to edit.
\end{quote}

\subsection*{You'll want to know: \ar{نشرة}}
\textbf{\ar{نشرة}} — \emph{nashra} — \textquotedblleft{}a bulletin, a news broadcast\textquotedblright{}.

\begin{quote}
\ar{شاهد} \ar{أبي} \ar{نشرة} \ar{الأخبار} \ar{في} \ar{المساء}.
\end{quote}

\begin{quote}
\emph{shāhada abī nashrat al-akhbār fī al-masāʾ.}
\end{quote}

\begin{quote}
My father watched the news bulletin in the evening.
\end{quote}

\ar{نشرة} \ar{الأخبار} is the news on radio or television.

\begin{cousinweb}
\textbf{\ar{نشرة}} (\emph{nashra}) is built on the native Arabic root \emph{n-sh-r} (Lane), \textquotedblleft{}to spread out, to unfold\textquotedblright{}, the root of the taught \emph{nashara}, to publish. A \emph{nashra} is what is spread among people.
\end{cousinweb}

\begin{grammarlens}[title={What you've built}]
One more word for telling what happened.
\end{grammarlens}

\subsection*{Your turn}
\begin{itemize}
\raggedright
  \item \emph{Say it:} \emph{nashra}
  \item \emph{Say it:} \emph{nashra}, once more
  \item \emph{Say it:} \emph{shāhada abī nashrat al-akhbār fī al-masāʾ.} — My father watched the news bulletin in the evening.
  \item \emph{Recall:} say the Arabic for an article, then the Arabic for to edit, then say \emph{nashra} again
\end{itemize}

\begin{tcolorbox}[breakable,colback=teal!4,colframe=teal!35!black,title={Before you move on}]
What does \ar{نشرة} mean? (\textquotedblleft{}A bulletin, a news broadcast\textquotedblright{}.) Say it once more.
\end{tcolorbox}
